\documentclass[10pt,a4paper]{article}
\usepackage{fontspec,amsmath,amssymb,amsthm,mathtools,hyperref}
\usepackage{thmtools,thm-restate,cleveref}
\newtheorem{theorem}{Theorem}
\newtheorem{lemma}[theorem]{Lemma}
\newtheorem{corollary}[theorem]{Corollary}
\newtheorem{definition}[theorem]{Definition}
\newtheorem{example}[theorem]{Example}
\newtheorem{remark}[theorem]{Remark}
\newtheorem{claim}[theorem]{Claim}
\newtheorem{proposition}[theorem]{Proposition}
\renewcommand{\epsilon}{\varepsilon}
\renewcommand{\theta}{\vartheta}
\usepackage{amsmath}
\usepackage{amssymb}
\usepackage{mathtools}

\usepackage{subcaption}
\usepackage{tikz}
\usetikzlibrary{math}
\usetikzlibrary{positioning,arrows.meta,fit,snakes,shapes}



\definecolor{rwth-blue}{cmyk}{1,.5,0,0}\colorlet{rwth-lblue}{rwth-blue!50}\colorlet{rwth-llblue}{rwth-blue!25}
\definecolor{rwth-violet}{cmyk}{.6,.6,0,0}\colorlet{rwth-lviolet}{rwth-violet!50}\colorlet{rwth-llviolet}{rwth-violet!25}
\definecolor{rwth-purple}{cmyk}{.7,1,.35,.15}\colorlet{rwth-lpurple}{rwth-purple!50}\colorlet{rwth-llpurple}{rwth-purple!25}
\definecolor{rwth-carmine}{cmyk}{.25,1,.7,.2}\colorlet{rwth-lcarmine}{rwth-carmine!50}\colorlet{rwth-llcarmine}{rwth-carmine!25}
\definecolor{rwth-red}{cmyk}{.15,1,1,0}\colorlet{rwth-lred}{rwth-red!50}\colorlet{rwth-llred}{rwth-red!25}
\definecolor{rwth-magenta}{cmyk}{0,1,.25,0}\colorlet{rwth-lmagenta}{rwth-magenta!50}\colorlet{rwth-llmagenta}{rwth-magenta!25}

\definecolor{rwth-yellow}{cmyk}{0,0,1,0}\colorlet{rwth-lyellow}{rwth-yellow!50}\colorlet{rwth-llyellow}{rwth-yellow!25}
\definecolor{rwth-grass}{cmyk}{.35,0,1,0}\colorlet{rwth-lgrass}{rwth-grass!50}\colorlet{rwth-llgrass}{rwth-grass!25}
\definecolor{rwth-green}{cmyk}{.7,0,1,0}\colorlet{rwth-lgreen}{rwth-green!50}\colorlet{rwth-llgreen}{rwth-green!25}
\definecolor{rwth-cyan}{cmyk}{1,0,.4,0}\colorlet{rwth-lcyan}{rwth-cyan!50}\colorlet{rwth-llcyan}{rwth-cyan!25}
\definecolor{rwth-teal}{cmyk}{1,.3,.5,.3}\colorlet{rwth-lteal}{rwth-teal!50}\colorlet{rwth-llteal}{rwth-teal!25}
\definecolor{rwth-gold}{cmyk}{.35,.46,.7,.35}
\definecolor{rwth-silver}{cmyk}{.39,.31,.32,.14}











\newenvironment{ansatz}[1][Ansatz]{\renewcommand{\qedsymbol}{$\lhd$}\begin{proof}[#1]}{\end{proof}}



\newcommand{\R}{\ensuremath{\mathbb{R}}}
\newcommand{\N}{\ensuremath{\mathbb{N}}}
\newcommand{\NN}{\ensuremath{\mathbb{N}}}
\newcommand{\Z}{\ensuremath{\mathbb{Z}}}
\newcommand{\Q}{\ensuremath{\mathbb{Q}}}

\newcommand{\A}{\ensuremath{\mathfrak A}}
\newcommand{\B}{\ensuremath{\mathfrak B}}
\newcommand{\C}{\ensuremath{\mathfrak C}}
\newcommand{\I}{\ensuremath{\mathfrak I}}
\newcommand{\LL}{\ensuremath{\mathfrak L}}

\newcommand{\TW}{\ensuremath{\mathcal{TW}}}
\newcommand{\TD}{\ensuremath{\mathcal{TD}}}
\newcommand{\ED}{\ensuremath{\mathcal{ED}}}
\newcommand{\E}{\ensuremath{\mathcal{T}}}
\newcommand{\Ekq}{\ensuremath{\mathcal{T}_{q}^{k}}}
\newcommand{\EParam}[2]{\ensuremath{\mathcal{T}_{#2}^{#1}}}
\newcommand{\Lkq}{\ensuremath{\mathcal{L}_{q}^{k}}}
\newcommand{\LParam}[2]{\ensuremath{\mathcal{L}_{#2}^{#1}}}
\newcommand{\GEkq}{\ensuremath{\mathcal{G}\Lkq}}
\newcommand{\GEParam}[2]{\ensuremath{\mathcal{G}\LParam{#1}{#2}}}
\newcommand{\GE}{\ensuremath{\mathcal{G}\mathcal{L}}}
\newcommand{\GEkqLL}{\ensuremath{\mathcal{G}\Ekq}}
\newcommand{\GEParamLL}[2]{\ensuremath{\mathcal{G}\LParam{#1}{#2}}}
\newcommand{\T}{\ensuremath{\mathcal{T}}}
\newcommand{\D}{\ensuremath{\mathcal D}}
\newcommand{\CG}{\ensuremath{\mathcal G}}

\newcommand{\lFO}{\ensuremath{\mathsf{FO}}}
\newcommand{\lC} {\ensuremath{\mathsf C}}
\newcommand{\lL} {\ensuremath{\mathsf L}}
\newcommand{\lGC}{\ensuremath{\mathsf{GC}}}

\newcommand{\xor}{\ensuremath{\oplus}}
\newcommand{\impl}{\ensuremath{\rightarrow}}

\newcommand{\parto}{\ensuremath{\rightharpoonup}}

\newcommand{\tup}[1]{\ensuremath{\boldsymbol{#1}}}
\newcommand{\qg}[1]{\ensuremath{\mathfrak{#1}}} \newcommand{\qgp}[3]{\ensuremath{\mathfrak{#1}[#3; #2]}} 

\DeclareMathOperator{\nonstatCR}{CR_{non-stat}}
\DeclareMathOperator{\CR}{CR}
\DeclareMathOperator{\monCR}{mon-CR}
\DeclareMathOperator{\supp}{supp}
\DeclareMathOperator{\qr}{qr}
\DeclareMathOperator{\lab}{lab}
\DeclareMathOperator{\free}{free}
\DeclareMathOperator{\vars}{vars}
\DeclareMathOperator{\dom}{dom}
\DeclareMathOperator{\img}{img}
\DeclareMathOperator{\HOM}{Hom}
\DeclareMathOperator{\pphom}{pp-hom}
\DeclareMathOperator{\pihom}{pi-hom}
\DeclareMathOperator{\tw}{tw}
\DeclareMathOperator{\td}{td}
\DeclareMathOperator{\wid}{wd}
\DeclareMathOperator{\dep}{dp}
\DeclareMathOperator{\ed}{ed}
\DeclareMathOperator{\BP}{BP}
\DeclareMathOperator{\cl}{cl}
\newcommand\restrict[1]{\vert_{#1}}
\DeclareMathOperator{\set}{set}

\renewcommand{\vec}[1]{\ensuremath{\boldsymbol{#1}}}

\newcommand{\elimOrd}[2]{$(#1,#2)$-constructible}
\newcommand{\ElimOrd}[2]{$(#1,#2)$-constructible}
\newcommand{\elimDepth}{elimination depth}
\newcommand{\labels}[1]{\ensuremath{L_{#1}}}
\newcommand{\grid}[2]{\ensuremath{G_{#1 \times #2}}}

\makeatletter
\let\amsmath@bigm\bigm

\renewcommand{\bigm}[1]{\ifcsname fenced@\string#1\endcsname
    \expandafter\@firstoftwo
  \else
    \expandafter\@secondoftwo
  \fi
  {\expandafter\amsmath@bigm\csname fenced@\string#1\endcsname}{\amsmath@bigm#1}}

\newcommand{\DeclareFence}[2]{\@namedef{fenced@\string#1}{#2}}
\makeatother

\DeclareFence{\mid}{|}

\newtheorem{fact}[definition]{Fact}
\Crefname{fact}{Fact}{Facts}



 
\usepackage[margin=22mm]{geometry}
\usepackage{url}
\providecommand{\doi}[1]{\href{https://doi.org/#1}{doi:#1}}
\newenvironment{claimproof}{\begin{proof}[Proof of claim]}{\end{proof}}
\providecommand{\claimqedhere}{\qedhere}

\usepackage{ifthen}

\definecolor{lipicsYellow}{rgb}{0.99,0.78,0.07}\colorlet{rwth-orange}{lipicsYellow}

\newtheorem{observation}[theorem]{Observation}

\usepackage{enumerate}
